
% ============================================================================
\clearpage
\section*{v4.8.0-r14 Source-Only: Separating Colour and Generation}
\addcontentsline{toc}{section}{v4.8.0-r14 Source-Only: Separating Colour and Generation}
\label{sec:v480-r14-colour-generation}
% ============================================================================

\paragraph{Purpose.}
Layer r13 recorded that Theorem~\ref{thm:three-generations} puts colour and generation on the same octonionic
triplet.  This layer examines two ways of separating them, both built from structures already present in the
Companion, and records the consequences for layers r12 and r13.  It is model construction; no prediction is made
and no data are compared.  Earlier layers are not rewritten.

\noindent\textbf{Source-only status.}
This v4.8.0-r14 layer is source-only.  It does not build a PDF, does not assign a new DOI,
and does not authorize a public release.  The v4.8.0 Version DOI
\href{https://doi.org/10.5281/zenodo.23365244}{10.5281/zenodo.23365244} and the Concept DOI
\href{https://doi.org/10.5281/zenodo.19210728}{10.5281/zenodo.19210728} are unchanged by it.

\subsection*{v4.8.0-r14 requirements}
\label{subsec:v480-r14-requirements}

A separation should (a) keep colour as \(SU(3)\subset G_2\), the stabilizer of \(e_7\); (b) count generations by a
symmetry that commutes with \(G_2\), so that a broken generation symmetry does not break colour; and (c) provide
nine linearly independent coloured states, three colours times three generations.

\subsection*{v4.8.0-r14 route A: sedenions}
\label{subsec:v480-r14-sedenions}

\paragraph{Theoremlet v4.8.0-r14.1.}
Extend the Cayley--Dickson convention of layer r04 once more, \(\mathbb S=\OO\oplus\OO e_8\).  The maps
\(\psi\), a rotation by \(2\pi/3\) in each of the seven planes \((e_i,e_{i+8})\), and \(\varepsilon\colon
a+be_8\mapsto a-be_8\) are automorphisms of \(\mathbb S\), satisfy \(\psi^3=\varepsilon^2=1\) and
\(\varepsilon\psi\varepsilon=\psi^{-1}\), and commute with \(G_2\) acting diagonally.  On the fourteen plane
directions \(1+\psi+\psi^2=0\).  Under \(G_2\times S_3\),
\[
\mathbb S=(\mathbf1,\mathbf1)\oplus(\mathbf1,\mathrm{sgn})\oplus(\mathbf7,\mathbf2).
\]
Coloured states therefore carry only the two-dimensional representation of \(S_3\): the nine vectors
\(\psi^k(e_c)\), \(c=4,5,6\), \(k=0,1,2\), span a space of dimension \(6\).

\paragraph{Proof sketch.}
Automorphism, group relations and commutation were checked numerically for random arguments, for all \(1344\)
signed-permutation automorphisms and for random elements \(\exp D\) of the fourteen-dimensional derivation
algebra (script \texttt{k28\_sedenion\_generationen.py}); the colour stabilizer of \(e_7\) has dimension \(8\).
The traces \(\operatorname{tr}\psi=-5\), \(\operatorname{tr}\varepsilon=0\) fix the decomposition.  That
\(\operatorname{Aut}(\mathbb S)\) is exactly \(G_2\times S_3\) is taken from the literature (Gillard and
Gresnigt, arXiv:2306.13098, and references there), not re-proved.  \(\square\)

Route A separates colour from generation, but the triole relation \(z_1+z_2+z_3=0\) itself makes the three
copies of each coloured state linearly dependent.  It yields only two independent coloured generations.

\subsection*{v4.8.0-r14 route B: the triality form}
\label{subsec:v480-r14-triality}

\paragraph{Theoremlet v4.8.0-r14.2.}
On \(\OO\oplus\OO\oplus\OO\) the trilinear form \(t(x,y,z)=\operatorname{Re}\bigl((xy)z\bigr)\) is invariant under the
cyclic shift \((x,y,z)\mapsto(y,z,x)\), under \((x,y,z)\mapsto(\bar x,\bar z,\bar y)\), which together generate
\(S_3\), and under \(G_2\) acting diagonally, which commutes with this \(S_3\).  The three copies of the colour
triplet \((e_4,e_5,e_6)\) give nine linearly independent coloured states.

\paragraph{Proof sketch.}
The real part of the associator vanishes and \(\operatorname{Re}(ab)=\operatorname{Re}(ba)\), so
\(\operatorname{Re}((xy)z)=\operatorname{Re}(x(yz))=\operatorname{Re}((yz)x)\).  For the reflection,
\(\operatorname{Re}((\bar x\bar z)\bar y)=\operatorname{Re}(y(zx))=\operatorname{Re}((xy)z)\).  Elements of
\(G_2\) preserve multiplication and conjugation.  Independence holds because the three copies lie in different
summands; all statements were also checked numerically.  \(\square\)

\paragraph{Proposed assignment.}
Generation index = the three triality slots; colour index = the \(SU(3)\) triplet \(e_4,e_5,e_6\) within each
slot.  This matches the Companion's identification of \(S_2\cong S^8=\OO P^1\) with three generations, and the
fact that \(G_2\) is the subgroup of \(\mathrm{Spin}(8)\) fixed by triality.  The assignment is proposed, not
derived; a native carrier for the three slots is still to be supplied.

\subsection*{v4.8.0-r14 consequences}
\label{subsec:v480-r14-consequences}

\begin{itemize}
\item After the separation, the order-\(96\) stabilizer of Theoremlet~r12.2 lies inside the colour group and
cannot act as a flavour symmetry.  Theoremlets r12.2--r12.4 remain correct mathematics but lose their reading
as lepton-mixing statements; H4 and H8 were already disfavoured by the JUNO 2026 re-check of layer r13.
\item A pure generation \(S_3\) with the triole \(\mathbb Z_3\) for charged leptons and a transposition for
neutrinos preserves the column \((0,\tfrac12,\tfrac12)\), i.e.\ \(\theta_{13}=0\), which is excluded.  This is
the minimality outcome of H6.  Realistic mixing requires further breaking, from which no prediction follows yet.
\item The triole form \(t(1,1,1)+\sqrt2\,\mathrm{Re}(c(1,\omega,\omega^2))\) of layer r10 is the decomposition of
generation space under \(S_3\) into the symmetric direction and the doublet; \(K=2/3\) remains the equal-norm
condition between them and is still not derived.
\end{itemize}

\subsection*{v4.8.0-r14 ledger and guard update}
\label{subsec:v480-r14-ledger-guard}

\[
R_{\rm cg}^{v480\text{-}r14}=(0,1,1)
\]
for (separation consistent with \(G_2\) achieved by route B; native carrier of the triality slots open; mixing
mechanism open).

The layer spends no guard:
\[
R_{\rm glob}^{(v480\text{-}r14,+)}=(0,1,1,1,1),
\qquad
\Gamma_{\rm guard}^{v480\text{-}r14}=(1,1,1,1,1),
\qquad
\beta_4=1/256\ \text{unspent}.
\]
It does not imply exact finite \(C5\) zero closure, exact all-readout terminality,
\(C6/T2\), a global photon-clock/RH theorem, or global all-readout closure.
